\documentclass[11pt]{article}
\usepackage[utf8]{inputenc}
\usepackage{amsmath,amssymb,amsthm}
\usepackage{booktabs}
\usepackage{geometry}
\usepackage{tikz}
\usetikzlibrary{fit,tikzmark,shapes.geometric}
\geometry{a4paper,margin=25mm}

\begin{document}

\section{Evaluation Methodology}

% ──────────────────────────────────────────────
% STEP 1 · LLM Stability
% ──────────────────────────────────────────────
\subsection{Step 1: LLM Stability Verification}
\label{sec:stability}

Before using LLM judgments as ground truth, we verify that the LLM produces \emph{consistent} results across repeated calls, regardless of correctness. We run each issue through the LLM five times independently and measure two complementary stability metrics.

\vskip 0.4em
\noindent\textbf{Set agreement } $S_{\mathrm{set}}$ measures the overlap of positive responses across runs:
\begin{equation}
S_{\mathrm{set}} \;=\; \frac{\bigl|\bigcap_{i=1}^{m} A_i\bigr|}
                           {\bigl|\bigcup_{i=1}^{m} A_i\bigr|}
\qquad
A_i = \{\,\text{issues labelled \emph{Yes} in run } i\,\}
\end{equation}
\vskip 0.3em
where $m=5$ is the number of runs. $S_{\mathrm{set}}=1$ means every run labels exactly the same set of issues as \emph{Yes}; lower values indicate increasing disagreement.

\vskip 0.4em
\noindent\textbf{Normalised entropy } $S_{\mathrm{entropy}}$ measures per-issue consistency:
\begin{equation}
S_{\mathrm{entropy}} \;=\; 1 - \bar{H}
\qquad
\bar{H} = \frac{1}{n}\sum_{j=1}^{n} H_j
\end{equation}
where for each issue $j$, $H_j = -p_j\log_2 p_j - (1-p_j)\log_2(1-p_j)$ is the binary entropy of its $m$ responses ($p_j$ = fraction of \emph{Yes} answers), and $n$ is the total number of issues. $S_{\mathrm{entropy}}=1$ means all issues receive identical answers across all runs.

\vskip 0.4em
Table~\ref{tab:stability} compares the two methods on the 200 retained issues, using five runs per method.

\begin{table}[htbp]
  \caption{LLM Stability Metrics for AgentBug-Smith (Ours) vs.\ Raw AI API ($m=5$ runs)}
  \label{tab:stability}
  \centering
  \begin{tabular}{lcc}
    \toprule
    \textbf{Method} & \textbf{$S_{\mathrm{set}}$} & \textbf{$S_{\mathrm{entropy}}$} \\
    \midrule
    AgentBug-Smith (Ours) & \textbf{0.9158} & \textbf{0.9378} \\
    Raw AI API            & 0.7208          & 0.8929          \\
    \bottomrule
  \end{tabular}
\end{table}

The \textbf{AgentBug-Smith (Ours)} method achieves $S_{\mathrm{set}}=0.9158$ and $S_{\mathrm{entropy}}=0.9378$, substantially outperforming the \textbf{Raw AI API} baseline ($S_{\mathrm{set}}=0.7208$, $S_{\mathrm{entropy}}=0.8929$). This confirms that the agent-issue taxonomy and structured prompt used by AgentBug-Smith produce markedly more consistent judgments across repeated LLM calls than a direct, unstructured prompt.

% ──────────────────────────────────────────────
% STEP 2 · Cochran's sampling formula
% ──────────────────────────────────────────────
\subsection{Step 2: Sampling Strategy}
\label{sec:sampling}

From the full dataset of $N=200$ labelled issues, we select $n=132$ for detailed human annotation. The sample size is determined by Cochran's formula for proportions:
\begin{equation}
n_0 \;=\; \frac{z^2 \; p(1-p)}{e^2}
\qquad
n \;=\; \frac{n_0}{1 + \dfrac{n_0-1}{N}}
\end{equation}
\vskip 0.3em
where $z=1.96$ (95\% confidence), $p=0.5$ (maximum variance assumption), $e=0.05$ (target margin of error), and $N=200$ (population size). This yields $n_0 \approx 384.2$, which after finite-population correction gives $n \approx 131.8$, rounded up to $n=132$.

% ──────────────────────────────────────────────
% STEP 3 · Human annotation + Kappa + Accuracy
% ──────────────────────────────────────────────
\subsection{Step 3: Human Annotation and Ground-Truth Construction}
\label{sec:annotation}

The 132 sampled issues are independently labelled by two annotators (the first author, \textbf{Annotator A}, and a domain expert, \textbf{Annotator B}) using the agent-issue taxonomy in \cite{agentbugsmith}. Each annotator classifies every issue as \emph{Yes} (agent issue) or \emph{No} (not an agent issue).

\vskip 0.4em
\noindent\textbf{Inter-annotator agreement.}  We measure agreement with Cohen's $\kappa$:
\begin{equation}
\kappa \;=\; \frac{p_o - p_e}{1 - p_e}
\end{equation}
where
\begin{equation}
p_o = \frac{1}{n}\sum_{i=1}^{n}\mathbf{1}[a_i = b_i]
\qquad
p_e = \sum_{c \in \{\mathrm{Y},\mathrm{N}\}} p_c^{(a)}\,p_c^{(b)}
\end{equation}
$p_o$ is observed agreement, $p_e$ is expected agreement by chance, $a_i$ and $b_i$ are the labels assigned by annotators $a$ (Annotator A) and $b$ (Annotator B) to issue $i$, and $p_c^{(a)}$, $p_c^{(b)}$ are the marginal proportions of class $c$ for each annotator.

\begin{table}[htbp]
  \caption{Human Annotation Statistics ($n=132$)}
  \label{tab:annotation}
  \centering
  \begin{tabular}{lc}
    \toprule
    \textbf{Metric} & \textbf{Value} \\
    \midrule
    Total labelled issues                         & 132   \\
    \midrule
    Annotator A labels --- Yes / No                   & 129 / 3   \\
    Annotator B labels --- Yes / No                    & 129 / 3   \\
    Human Golden --- Yes / No                    & 128 / 4   \\
    \midrule
    Observed agreement $p_o$                     & $128/132 = \mathbf{0.9697}$ \\
    Chance agreement $p_e$                       & $0.9654\times0.9654 + 0.0346\times0.0346 = \mathbf{0.9556}$ \\
    Cohen's $\kappa$                             & $\mathbf{0.3178}$ \\
    \midrule
    Annotator A accuracy vs.\ Human Golden            & $131/132 = \mathbf{0.9924}$ \\
    Annotator B accuracy vs.\ Human Golden             & $129/132 = \mathbf{0.9773}$ \\
    \bottomrule
  \end{tabular}
\end{table}

\noindent The annotator contingency table is
\[
\begin{array}{c|cc}
 & \text{Annotator B: Yes} & \text{Annotator B: No} \\
\hline
\text{Annotator A: Yes} & 127 & 2 \\
\text{Annotator A: No}  & 2   & 1
\end{array}.
\]
The off-diagonal counts are balanced (2 vs.\ 2), indicating no directional tendency for one annotator to assign more positive labels than the other. Positive agreement is high,
\[
P_{\mathrm{pos}}
=\frac{2(127)}{2(127)+2+2}
=0.9845,
\]
whereas negative agreement is
\[
P_{\mathrm{neg}}
=\frac{2(1)}{2(1)+2+2}
=0.3333.
\]
Thus, the annotators agree strongly on positive cases, while the small number of negative cases is intrinsically less stable.

\vskip 0.4em
\noindent\textbf{Adjudication.} After independent annotation, the disagreements are reviewed and Annotator B assigns the final adjudicated label, denoted \textbf{Human Golden}. This label is used as the reference for downstream validation. Because Annotator B performs the final adjudication, the reported agreement between Annotator B's initial labels and Human Golden is descriptive rather than an independent accuracy estimate.

\vskip 0.4em
\noindent\textbf{Interpretation standard for $\kappa$.} We use the commonly cited Landis--Koch interpretation ranges~\cite{landis1977measurement}:
\begin{table}[htbp]
  \caption{Landis--Koch Interpretation of Cohen's $\kappa$}
  \label{tab:kappa-ranges}
  \centering
  \begin{tabular}{lc}
    \toprule
    \textbf{Cohen's $\kappa$ range} & \textbf{Strength of agreement} \\
    \midrule
    $\kappa < 0.00$       & Poor \\
    $0.00$--$0.20$        & Slight \\
    \tikzmarknode{kappal}{$0.21$--$0.40$} & \tikzmarknode{kappar}{Fair} \\
    $0.41$--$0.60$        & Moderate \\
    $0.61$--$0.80$        & Substantial \\
    $0.81$--$1.00$        & Almost perfect \\
    \bottomrule
  \end{tabular}
  \begin{tikzpicture}[overlay,remember picture]
    \node[draw=red,very thick,inner xsep=6pt,inner ysep=2pt,fit=(kappal)(kappar)] {};
  \end{tikzpicture}
\end{table}

\noindent Our result, $\kappa=0.318$, falls within the \emph{fair agreement} range ($0.21$--$0.40$). However, the raw agreement $p_o=96.97\%$ is very high. This apparent paradox is a known artefact of Cohen's $\kappa$ under severe class imbalance: when one class dominates (129 of 132 initial labels are \emph{Yes} for each annotator), even a small number of disagreements sharply reduces $\kappa$~\cite{feinstein1990high}. Therefore, the categorical range should be treated as descriptive rather than as an absolute assessment of annotation quality. As a prevalence-robust sensitivity check, Gwet's AC1 is $0.9683$, which is consistent with the observed agreement and confirms that the modest Cohen's $\kappa$ is primarily caused by prevalence rather than broad annotator disagreement.

% ──────────────────────────────────────────────
% STEP 4 · LLM Accuracy vs. Human Ground Truth
% ──────────────────────────────────────────────
\subsection{Step 4: Ours vs.\ Raw LLM against Human Golden}
\label{sec:accuracy}

Human Golden on the $n=132$ sample comprises 128 \emph{Yes} and 4 \emph{No}. We compare two predictors against this reference.

\begin{itemize}
  \item \textbf{Ours} (AgentBug-Smith): the taxonomy-guided filter; every sampled issue already received \emph{Yes}.
  \item \textbf{Raw LLM}: \texttt{gpt-4.1-mini}, five independent votes on issue title and body only (no patch, no \texttt{ai\_judgment}, no taxonomy), aggregated by majority.
\end{itemize}

\noindent\textbf{Confusion counts.}
\[
\begin{array}{rcl}
\mathrm{TP} &=& \#(\mathrm{pred}=\mathrm{Yes},\;\mathrm{Golden}=\mathrm{Yes}) \\
\mathrm{FP} &=& \#(\mathrm{pred}=\mathrm{Yes},\;\mathrm{Golden}=\mathrm{No}) \\
\mathrm{TN} &=& \#(\mathrm{pred}=\mathrm{No},\;\mathrm{Golden}=\mathrm{No}) \\
\mathrm{FN} &=& \#(\mathrm{pred}=\mathrm{No},\;\mathrm{Golden}=\mathrm{Yes})
\end{array}
\qquad
\begin{array}{rcl}
\mathrm{Acc} &=& (\mathrm{TP}+\mathrm{TN})/n \\
\mathrm{P} &=& \mathrm{TP}/(\mathrm{TP}+\mathrm{FP}) \\
\mathrm{R} &=& \mathrm{TP}/(\mathrm{TP}+\mathrm{FN}) \\
\mathrm{Spec} &=& \mathrm{TN}/(\mathrm{TN}+\mathrm{FP}) \\
F_1 &=& 2\mathrm{PR}/(\mathrm{P}+\mathrm{R})
\end{array}
\]
Because Golden has only four negatives, $\mathrm{FP}+\mathrm{TN}=4$ and $\mathrm{TP}+\mathrm{FN}=128$ for both methods.

\begin{table}[htbp]
  \caption{Ours vs.\ Raw LLM: confusion \emph{counts} against Human Golden ($n=132$)}
  \label{tab:ours-vs-raw-counts}
  \centering
  \begin{tabular}{lcccc}
    \toprule
    \textbf{Method} & TP & FP & TN & FN \\
    \midrule
    Ours     & \textbf{128} & 4 & 0 & \textbf{0} \\
    Raw LLM  & 70  & 4 & 0 & 58 \\
    \bottomrule
  \end{tabular}
\end{table}

\begin{table}[htbp]
  \caption{Ours vs.\ Raw LLM: cell percentages of $n=132$ ($\mathrm{TP}\%=\mathrm{TP}/n$, etc.)}
  \label{tab:ours-vs-raw-cellpct}
  \centering
  \begin{tabular}{lcccc}
    \toprule
    \textbf{Method} & TP\% & FP\% & TN\% & FN\% \\
    \midrule
    Ours     & $\mathbf{128/132=96.97\%}$ & $4/132=3.03\%$ & $0\%$ & $\mathbf{0\%}$ \\
    Raw LLM  & $70/132=53.03\%$ & $4/132=3.03\%$ & $0\%$ & $58/132=43.94\%$ \\
    \bottomrule
  \end{tabular}
\end{table}

\begin{table}[htbp]
  \caption{Ours vs.\ Raw LLM: confusion \emph{rates} against Human Golden}
  \label{tab:ours-vs-raw-rates}
  \centering
  \small
  \begin{tabular}{lccccc}
    \toprule
    \textbf{Method} & Acc & Precision & Recall & Spec & $F_1$ \\
    \midrule
    Ours     & $\mathbf{128/132=96.97\%}$ & $\mathbf{128/132=96.97\%}$ & $\mathbf{128/128=100\%}$ & $0/4=0\%$ & $\mathbf{98.46\%}$ \\
    Raw LLM  & $70/132=53.03\%$ & $70/74=94.59\%$ & $70/128=54.69\%$ & $0/4=0\%$ & $69.31\%$ \\
    \bottomrule
  \end{tabular}
\end{table}

Both methods label all four Golden-\emph{No} issues as \emph{Yes} ($\mathrm{FP}=4$, $\mathrm{TN}=0$), so specificity is 0 in either case. The gap is recall: Ours never misses a Golden-\emph{Yes} issue ($\mathrm{FN}=0$, $\mathrm{R}=100\%$), whereas Raw LLM misses $58/128=54.69\%$ of them. Accuracy therefore falls from $96.97\%$ (Ours) to $53.03\%$ (Raw LLM). Ours' recall of $100\%$ holds only on this already-filtered sample, where the method predicts \emph{Yes} for every case.

% ──────────────────────────────────────────────
% References
% ──────────────────────────────────────────────
\begin{thebibliography}{99}
  \bibitem{agentbugsmith} Eamin et al., \textit{AgentBug-Smith: A Taxonomy of Agent Issues in LLM-based Agent Systems}, 2025.
  \bibitem{landis1977measurement} Landis, J. R., and Koch, G. G., ``The measurement of observer agreement for categorical data,'' \textit{Biometrics}, vol. 33, no. 1, pp. 159--174, 1977.
  \bibitem{feinstein1990high} Feinstein, A. R., and Cicchetti, D. V., ``High agreement but low kappa: II. Resolving the paradoxes,'' \textit{Journal of Clinical Epidemiology}, vol. 43, no. 6, pp. 551--558, 1990.
\end{thebibliography}

\end{document}
